\chapter{The Proxy-Bias Result}
\label{ch:proxybias}

Black-box prompt optimization is only affordable because candidates are scored on a small
subset of the benchmark rather than the whole. This chapter shows that the choice of subset
is not a harmless implementation detail: the intuitive ``keep the most informative items''
rule is a biased estimator of the full-benchmark score, the bias cannot be removed by
calibration, and---as Chapter~\ref{ch:repair} shows---it can steer the entire optimization
to the wrong answer.

\section{The symptom}

The first version of our pipeline selected the 180 proxy items whose binary outcomes varied
most across prompt configurations (``discrimination'' selection), on the standard intuition
that the most informative items should best track the benchmark. When the optimizer's chosen
configurations were validated on the full benchmark, their scores came in systematically
below their proxy scores---by 8 to 26 MB points, and by an amount that differed per
configuration. A join of proxy-predicted against full scores gave a high Pearson correlation
($r \approx 0.95$) but this was computed on only four points (one per model); the same four
points give a Spearman rank correlation of $1.0$ on MB and $0.8$ on HA. Correlation on
$n=4$ is not evidence, and the offline study below shows the discrimination proxy is in fact
the \emph{worst} ranker of all the strategies we tested, contradicting the very property it
was chosen for.

\section{The offline leave-one-prompt-out study}

Because the pipeline had saved every model's per-item answer for 14 prompt configurations
across 4 models (56 configurations in total), we could measure the quality of any selection
strategy by \emph{replay}, at zero API cost. The protocol is leave-one-prompt-out, per
model: to predict a held-out prompt's full-benchmark HA and CP, the proxy items are selected
using only the \emph{other} prompts, and the held-out prompt's score is then estimated from
those items. This is an out-of-sample estimate, not a fit.

A subtlety in the original analysis inflated the apparent advantage of representative
sampling: predictions were averaged over five random draws \emph{before} the error was
computed. Five draws of 90 items from a pool of ${\sim}357$ cover most of the pool, so the
averaged number describes a near-complete evaluation rather than a single deployable proxy.
We therefore report the honest \emph{single-draw} error, with the ensembled number kept only
as a diagnostic. Table~\ref{tab:lopo} gives the corrected result.

\begin{table}[H]
\centering
\caption{Single-draw leave-one-prompt-out error (180-item proxy, 56 configurations, 95\% CI).}
\label{tab:lopo}
\begin{tabular}{lcccc}
\toprule
selection & MB MAE [95\% CI] & HA MAE & CP MAE & MB $r$ \\
\midrule
variance / discrimination (original) & 12.9 [9.9, 16.0] & 11.0 & 8.5 & 0.81 \\
difficulty-stratified & 4.2 [3.3, 5.3] & 3.6 & 3.1 & 0.98 \\
random & 2.5 [2.0, 2.9] & 1.8 & 1.4 & 0.999 \\
\bottomrule
\end{tabular}
\end{table}

The discrimination proxy is the worst on every metric, \emph{including ranking} ($r$ 0.81 vs
0.999). The gap to representative sampling is roughly three-to-fivefold---not the ``order of
magnitude'' an earlier draft claimed once the unfair ensembling is removed. But the raw ratio
understates the real distinction, which is one of \emph{kind}: the discrimination proxy's
${\sim}13$-point error is \textbf{bias} (identical under single-draw and five-draw
scoring---averaging does not reduce it), whereas the representative samplers' 2--4-point error
is \textbf{sampling noise} (random's error falls from 2.5 to 0.8 when five draws are
ensembled). Bias does not shrink with more measurement; noise does. This is why
representative sampling is the correct fix and no amount of re-running the biased proxy would
have helped.

\section{Why informativeness biases the mean}

The mechanism is a classical result from active testing \citep{kossen2021activetesting}. A
binary item's variance across configurations is maximized when its success probability is
near $0.5$, so a variance-maximizing rule preferentially keeps \emph{borderline} items.
Borderline items have outcome means near 50\% almost by definition; restricting the subset to
them pulls the subset mean toward 50\% and away from the true population mean. The size of the
distortion depends on how far the configuration's true score sits from 50\%, which differs per
configuration---hence the bias is configuration-conditional, and a single global correction
cannot remove it. A representative sample, by contrast, keeps easy and borderline items in
their true proportions, so its mean is an unbiased estimate with only sampling noise.

\section{The bias is a property of informativeness, not one heuristic}

Two further analyses (both zero-cost replay) sharpen the finding and connect it to current
practice.

\paragraph{Selection-principle ablation.} Recent performance-guided selection methods for
prompt optimization \citep{ipomp2025} combine a \emph{coverage} principle with a
\emph{boundary} principle (prefer decision-boundary / high-uncertainty items). Ablating the
two in our harness (Table~\ref{tab:ablation}) shows coverage is benign while boundary is
toxic: boundary selection alone reproduces the discrimination bias, and adding a boundary
term to an otherwise unbiased coverage-clustering selection triples its error
($3.0 \to 10.9$).

\begin{table}[H]
\centering
\caption{Selection-principle ablation (single-draw MB MAE).}
\label{tab:ablation}
\begin{tabular}{lc}
\toprule
principle & MB MAE \\
\midrule
random & 2.5 \\
coverage-clustering & 3.0 \\
difficulty-stratified (shipped) & 4.2 \\
coverage + boundary & 10.9 \\
variance / discrimination & 12.9 \\
boundary alone & 13.2 \\
\bottomrule
\end{tabular}
\end{table}

The failure is intrinsic to \emph{informativeness as a selection objective for score
estimation}, not to any particular heuristic: methods that select uncertain or boundary
items inherit the bias whenever the subset is also used to estimate an aggregate score.

\paragraph{No calibration fix.} Fitting $\text{proxy} = \alpha\cdot\text{full} +
(1-\alpha)\cdot 50$ gives $\alpha \approx 1.0$ for every strategy---the bias does not manifest
as a global affine attenuation, because its offset depends on the unknown true score of the
configuration being predicted. A per-model affine calibration reduces the discrimination
proxy's MB MAE only from 12.9 to 9.6, still three to four times worse than plain random
sampling. The bias must be prevented in selection, not repaired in scoring.

\section{A principled estimator does not help at this scale}

For completeness we implemented the tinyBenchmarks/metabench approach
\citep{polo2024tinybenchmarks,metabench2024}: a 2-parameter Item Response Theory model fit to
the configuration $\times$ item matrix, with anchor items chosen by quantile-stratified
Fisher information. At our scale the IRT estimator is worse than plain stratified or random
sampling at every proxy size: fitting hundreds of item parameters from only 56 configurations
is under-determined. Below that data threshold, the mean of a representative sample is hard to
beat, consistent with \citet{missesthemark2025}.

\section{Summary}

Selecting benchmark items for informativeness biases the estimated aggregate score in a
configuration-dependent way that no calibration can undo; representative sampling of the same
budget is nearly unbiased; and a sophisticated IRT estimator does not help at our data scale.
The next chapter shows that this is not merely a reporting problem: the same bias, inside an
optimization loop, changes which prompt the optimizer returns.
